\subsection{FRM Exam 1999---Question 89 (Market Risk): interpretation of VaR}

\textbf{Enunciado}

\emph{What is the correct
interpretation of a \$3 million overnight VAR figure with 99\% confidence
level? The institution}
\begin{enumerate}[label=\alph*.]
\item {Can be expected to lose at most \$3 million in 1 out of next 100 days}
\item {Can be expected to lose at least \$3 million in 95 out of next 100 days}
\item \hl{Can be expected to lose at least \$3 million in 1 out of next 100 days}
\item {Can be expected to lose at most \$6 million in 2 out of next 100 days}
\end{enumerate}

\subsubsection*{Solución}



Sea \(L\) la pérdida overnight. El VaR al nivel \(\alpha\) es el cuantil \(\alpha\):
\[
\VaR_\alpha(L):=\inf\{\ell:\Prob(L\le \ell)\ge \alpha\}.
\]
Que se escribe como:
\[
\Prob(L>\VaR_\alpha(L))\approx 1-\alpha.
\]

\textbf{Sustituir \(\alpha=0.99\).} Si \(\VaR_{0.99}(L)=\$3{,}000{,}000\), entonces:
\[
\Prob(L>3{,}000{,}000)\approx 0.01.
\]
Es decir: \emph{aprox. 1 de cada 100 días la pérdida excede (o es al menos) \$3 millones.}


